\section{System 2: Geometric Impedance and Causal Propagation Probability}
\label{sec:system2_geometric_impedance}

\CatchFileBetweenTags{\AlphaInvCAPVal}{constants.tex}{AlphaInvCAPVal}
\CatchFileBetweenTags{\AlphaInvMAPVal}{constants.tex}{AlphaInvMAPVal}
\CatchFileBetweenTags{\AlphaInvPROVal}{constants.tex}{AlphaInvPROVal}
\CatchFileBetweenTags{\AlphaInvGOVVal}{constants.tex}{AlphaInvGOVVal}
\CatchFileBetweenTags{\AlphaInvTOLVal}{constants.tex}{AlphaInvTOLVal}
\CatchFileBetweenTags{\AlphaInvMARVal}{constants.tex}{AlphaInvMARVal}
\CatchFileBetweenTags{\AlphaInvVal}{constants.tex}{AlphaInvVal}

The $E_8$ substrate is a maximally constrained, zero-flux causal network. At Quiescent Equilibrium there is no kinetic motion in the traditional sense. Yet persistence across discrete time requires that the substrate's structural state be transferred from one temporal node to the next. This is not dynamical propagation but \textbf{causal serialization}: the defect must re-assert its boundary at each clock tick or dissolve into the homogeneous background. In this discrete setting, the natural measure of structural survival is the network \textbf{admittance} $Y = 1/Z$.

On this geometric substrate, the admittance represents the fundamental \textbf{causal coupling} ($\kappa$): the probability that information successfully transfers to the next temporal node without decohering:
\begin{equation}
    \kappa \equiv Y = \frac{1}{Z_{\mathrm{geo}}}.
\end{equation}

On a homogeneous network, any localized state naturally tends to disperse into geometric noise. To remain persistently observable, a configuration must maintain a stable topological partition across its fundamental causal cycle. Under the mechanics of persistence, the structural cost of survival is governed by a variational impedance. The six orthogonal pillars act as competing geometric forces. They contribute either Open-Loop Structural Costs or Active Feedback Reductions, which sum linearly to the total geometric impedance.

% ============================================================
% CAP must satisfy all of the following:
%
% 1. BLINDING TEST. Every step must be one we would accept without
%    knowing alpha^-1 ~ 137. Any step persuasive only because the
%    total comes out right is numerology and not acceptable.
% 2. DEPENDENCY DIRECTION. Premises may only be System 0 requirements
%    (Closure, Finiteness, Conservation, Causality, Homogeneity) and
%    System 1 invariants (Delta=43, c=1, Euclidean substrate metric,
%    periodic causal cycle). Nothing from System 2 onward and no
%    companion-paper dynamics (Wilson loops, gravity strain,
%    Zitterbewegung) may be premises; they may appear only as labeled
%    forward-consistency remarks.
% 3. FORWARD CONSISTENCY. The result may not contradict later
%    structure (Wilson perimeter law, capacity-strain gravity, mass
%    as closed chiral circulation), but consistency is a CHECK, never
%    a premise.
% 4. PI MUST BE CONSTRUCTED. Pi may enter only as the measured
%    invariant of an explicitly constructed geometric object. The
%    microscopic lattice has no continuous rotational symmetry; the
%    isotropic round ball is the macroscopic idealization forced by
%    the Homogeneity requirement (macroscopic observables must be
%    invariant under lattice orientation). The step where isotropy
%    replaces the discrete lattice ball with the round ball must be
%    explicit, named, and scoped to macroscopic entropic cost.
% 5. EVERY FACTOR INTERROGATED. The text must answer, in-line:
%    why the factor is exactly 1; why diameter Delta and not radius
%    Delta (why not 2*pi*Delta); why the perimeter and not the area
%    (why not ~pi*Delta^2/4) or the path length (why not Delta); why
%    the sign is structural (+) under the Entropic Balance Sheet.
% 6. OPERATIONAL UNITS. The result must be a dimensionless count of
%    elementary link actions per fundamental cycle, consistent with
%    S_Phi as a transition count.
% 7. ONE-SENTENCE TEST. The geometric picture must be stateable in
%    one sentence: a persistent defect is a closed circulation whose
%    parts must remain in causal contact within one cycle; its
%    boundary must be re-traced every cycle.
% ============================================================
\subsection{The Resonant Circumference (Capacity)}
\label{subsec:alpha_cap}

% What happens if we dont' have CAP
Capacity contributes an \textbf{Open-Loop Structural Cost} that prevents \emph{instantaneous delocalization}. It is the absolute geometric allocation of the spatial envelope: the continuous perimeter required to enclose the configuration across its fundamental cycle. Because the defect is defined strictly by its topological boundary and its interior is indistinguishable from the vacuum, the entropic cost is a boundary measure, not a bulk volume.

% Derive the maximal reach from periodicity
At the network speed limit ($c=1$), a signal can traverse at most one lattice spacing per cycle. To maintain a single coherent boundary over the full $\Delta$ cycle, a signal must complete a round trip from the defect’s center to its edge and back, bounding the causal radius to $\Delta/2$ and the diameter to $D = \Delta$.

% Establish the metric and evaluation domain
At Quiescent Equilibrium, the substrate's footprint is evaluated on its native positive-definite (Euclidean) metric required by Conservation (\Cref{subsec:positive_definite}), with causal propagation restricted to a $(1+1)D$ plane $(x,\tau)$.

% circumference
By Homogeneity, macroscopic observables must be invariant under lattice orientation; any discrete anisotropy cancels in the orientation average. The invariant envelope (the union of all valid closed $\Delta$-cycles, averaged over all orientations) is therefore a continuous, isotropic disk of circumference $\pi D$ evaluated on the native Euclidean metric. Because the defect interior is indistinguishable from the vacuum, the entropic cost is the measure of its smooth boundary.

\begin{equation}
    Z_{\mathrm{CAP}} = C = \pi D = \pi\Delta \approx \num{\AlphaInvCAPVal}
\end{equation}

The dominant contribution to the vacuum impedance ($\approx 99.9\%$) originates from this fundamental geometry (this corresponds in the standard effective field theory dual to the Wilson Loop).

\subsection{Topological Boundary (Map)}
\label{subsec:alpha_map}
Map contributes an \textbf{Open-Loop Structural Cost} that prevents \emph{boundary dissolution}. It is the absolute allocation of the topological boundary ($\chi = 2$) that prevents internal state data from leaking into the unquantized background.

As derived in \Cref{sec:topological_boundary}, the minimal persistent defect is the simply-connected sphere ($S^2$), which carries the maximum Euler characteristic for any closed surface:
\begin{equation}
    Z_{MAP} = +\chi = \AlphaInvMAPVal
\end{equation}
Because $\chi$ is a pure topological invariant, it is exact and structurally forbidden from scaling, running, or acquiring continuous corrections.


\subsection{Alignment Efficiency (Protocol)}
\label{subsec:alpha_pro}
Protocol acts as an \textbf{Active Feedback Reduction} that prevents \emph{internal phase decoherence}. Unlike the absolute allocations of Capacity and Map, Protocol coordinates phase information across existing orthogonal channels (Sectors A and B). Its contribution is therefore a fractional ratio of the residual coordination capacity.

The manifold resolution $R_M = D\Delta$ (\Cref{subsec:manifold_resolution}) is the total available coordination capacity, the maximum distinct spacetime addresses per fundamental cycle. A localized defect consumes exactly $\sigma = 5$ degrees of freedom to define its geometric embedding, leaving residual capacity:

\begin{equation}
    C_{res} = R_M - \sigma = (4 \cdot 43) - 5 = 167
\end{equation}

Entropic impedance is the inverse of flow capacity. The greater the residual capacity, the lower the resistance to phase transport, and because efficient coordination reduces net dissipation, this contribution is an organizational gain, strictly negative in the overall structural budget:
\begin{equation}
    Z_{PRO} = -\frac{1}{C_{res}} = -\frac{1}{167} \approx \num{\AlphaInvPROVal}
\end{equation}


\subsection{Stabilizing Potential (Governor)}
\label{subsec:alpha_gov}
Governor acts as an \textbf{Active Feedback Reduction} that prevents \emph{feedback runaway}. It is the uniform distribution of the topological boundary constraint across the fundamental causal loop, converting the static boundary invariant into a stabilizing density that suppresses unbounded divergence.

By Homogeneity, this constraint cannot be anchored to any preferred temporal coordinate; it must be distributed maximally uniformly across the loop of length $\Delta = 43$. To prevent non-linear self-interactions which would violate the vacuum's minimal interaction rule (the Toll constraint), the geometric coupling must be strictly first-order (linear). The exact stabilizing potential density is therefore the linear ratio of the boundary ($\chi = 2$) to the loop length ($\Delta = 43$). As an organizational gain, its contribution is negative:
\begin{equation}
    Z_{GOV} = -\frac{\chi}{\Delta} = -\frac{2}{43} \approx \num{\AlphaInvGOVVal}
\end{equation}


\subsection{Entropic Transition (Toll)}
\label{subsec:alpha_tol}
Toll contributes an \textbf{Open-Loop Structural Cost} that prevents \emph{causal collapse} (Zeno breakdown). It is the eliminative filter that selects valid state transitions from the unconstrained combinatorial background. Its contribution is therefore a probability ratio: the minimal valid phase space divided by the maximal unconstrained phase space.

\begin{equation}
    Z_{TOL} = \frac{\Omega_{valid}}{\Omega_{max}}.
\end{equation}

The numerator, the minimal state transition, is defined by three geometric filters required to prevent non-unitary mixing, unquantized flux, or coordinate self-aliasing:
\begin{enumerate}
    \item \textbf{State Identity (1):} The interaction acts as the identity on internal channel configuration.
    \item \textbf{Topological Flux ($\chi=2$):} The boundary permits exactly two independent flux orientations.
    \item \textbf{Coordinate Non-Degeneracy ($R_M - \sigma$):} The transition target cannot alias with the interaction's own geometric embedding, leaving the orthogonal residual $(R_M - \sigma)$ valid coordinates.
\end{enumerate}

These three filters are independently necessary and jointly sufficient; no fourth filter is available because the $E_8\to D_4\oplus D_4$ projection exhausts the node's orthogonal geometric generators. Because identity (Sector B), topological flux (Sector A), and coordinate non-degeneracy ($R_M - \sigma$) govern strictly orthogonal configuration spaces, their independent valid state counts combine multiplicatively. The minimal valid phase space is therefore exactly:

\begin{equation}
    \Omega_{valid} = 1 \cdot \chi \cdot (R_M - \sigma).
\end{equation}

The denominator counts independent dimensions of the unconstrained configuration manifold. On the substrate, a non-trivial state transition requires an interaction; a 2-point operation is mere propagation (identity). To conserve information across distinct channel types, the lattice must couple them in symmetric groups of three: any fewer channels cannot mediate a non-trivial interaction, while any additional channel must duplicate an existing type and collapse distinct states into a single temporal slot. This geometric requirement forces the minimal interaction vertex to be a triality singlet: $\mathbf{8}_v \otimes \mathbf{8}_s \otimes \mathbf{8}_c \to \mathbf{1}$. Because the $D_4$ spinor structure provides exactly three independent 8-dimensional representations, a fourth leg would produce Causal Aliasing, rendering the causal history non-injective and violating Unitarity. Higher-point transitions are therefore composite sequences of 3-point vertices.

With no external observer to assign internal channel roles prior to interaction, an unconstrained 3-point vertex cannot pre-distinguish between Sector A and Sector B degrees of freedom. Rather than being restricted to the single-sector internal bit-depth $\nu=16$, each leg independently samples the total topological node capacity $N=2\nu=32$ across both orthogonal sectors.

The unconstrained phase space of admissible state assignments for a 3-leg vertex is therefore $N^3$. Combining this with the local interaction degrees of freedom ($\sigma=5$) and the available manifold addresses ($R_M=172$), the maximal unconstrained phase space is the Cartesian product:

\begin{equation}
    \Omega_{max} = N^3 \cdot \sigma \cdot R_M.
\end{equation}

The Toll is therefore the ratio of minimal valid to maximal unconstrained phase space:
\begin{equation}
\begin{split}
    Z_{TOL} &= \frac{\Omega_{valid}}{\Omega_{max}} 
    = \frac{\chi(R_M - \sigma)}{N^3 \sigma R_M} \\
    &\approx \num{\AlphaInvTOLVal}
\end{split}
\end{equation}


\subsection{Resolution Floor (Margin)}
\label{subsec:alpha_mar}
Margin contributes an \textbf{Open-Loop Structural Cost} that prevents \emph{stochastic fragility}. It is the resolution threshold that buffers the signal against geometric noise. Its contribution is therefore a fractional dilution factor: one quantum of structured information divided by the minimal configuration volume.

The minimum resolvable signal is exactly 1 quantum of information diluted across the \emph{minimal} number of distinct configurations ($V_{config}$) required to specify a localized defect. The minimal configuration volume requires three domains to exhaust the node's orthogonal geometric generators and prevent underspecification:
\begin{enumerate}
    \item \textbf{The State Vector ($L_{embed} = 31$):} The total description length of a dynamic defect, structurally derived in \cref{par:lembed}.
    \item \textbf{The Local Interaction Volume ($V_{local} = \sigma + 1 = 6$):} A localized lattice interaction is a connected subgraph spanning $\sigma = 5$ independent directions from a central node. To strictly localize the interaction to a single neighborhood, the maximum distance from the center to any node must be 1 (graph radius $r=1$). The Star Graph ($S_\sigma$) is the strictly unique connected tree topology of $\sigma$ edges that possesses a radius of 1; any other connected tree has radius $r \ge 2$, delocalizing the interaction across multiple neighborhoods and violating the Orthogonality Condition. The minimal local volume is therefore $\sigma+1=6$.
    \item \textbf{The Causal Area} ($A_{causal} = \Delta^2$): As established in \cref{sec:fundamental_resonance_filter1}, fundamental causality restricts information propagation on the discrete substrate to a local $(1+1)$D plane. Operating at the maximum speed limit $c \equiv 1$, a fundamental cycle of temporal duration $\Delta$ sweeps a 1D spatial horizon of length $\Delta$. The resulting spacetime resolution count is strictly $\Delta \cdot \Delta = \Delta^2$.
\end{enumerate}

Because state vector, local topology, and causal horizon constitute orthogonal configuration spaces, their independent specification domains combine via Cartesian product, yielding a multiplicative volume:
\begin{equation}
    V_{config} = L_{embed} \cdot (\sigma + 1) \cdot \Delta^2
\end{equation}

The Margin impedance is the reciprocal dilution factor: the fraction of the total configuration volume occupied by a single quantum of structured information:

\begin{equation}
\begin{split}
    Z_{MAR} &= \frac{1}{V_{config}} = \frac{1}{L_{embed} \cdot (\sigma + 1) \cdot \Delta^2} \\
    &=\frac{1}{31 \cdot (5+1) \cdot 43^2} \approx \num{\AlphaInvMARVal}
\end{split}
\end{equation}


\subsection{The Additive Impedance Theorem}
\label{subsec:additive_theorem}

\noindent\textbf{The Additive Impedance Theorem.} 
The total geometric impedance $Z_{\mathrm{geo}}$ is the exact linear combination of the six pillar impedances with unit coefficients. This additive form is not an approximation; it is the unique mathematical structure consistent with the existential requirements of System 0.

\vspace{0.5em}
\noindent\textit{Proof.} The theorem rests on four structural deductions forced by the Closure Constraint and the Quiescent Equilibrium:

\subsubsection{Universal Currency (Dimensional Homogeneity)}
Standard physics requires conversion constants ($c, \hbar, G$) to map heterogeneous phenomena into macroscopic unit conventions. The Closure Constraint forbids external reference frames, reducing all geometric measures to a single ontological currency: the dimensionless count of elementary state transitions. With the lattice spacing and fundamental clock tick structurally defined as unity ($a=1, \delta t=1$), spatial perimeter, topological boundaries, and state-space allocations map bijectively to state transitions. As counts of the identical primitive operation, their linear combination is mathematically well-defined without dimensional conversion factors.

\subsubsection{Unit Weighting and Kinematic Normalization}
In standard physics, fractional coefficients ($1/2$, $1/\sqrt{2}$) arise from external thermodynamic weights or internal kinematic projection norms. Here, the Closure Constraint forbids external thermal baths (Lagrange multipliers). Internally, because the $E_8$ substrate is strictly unimodular ($\text{vol}(\Lambda)=1$) and self-dual ($\Lambda = \Lambda^*$), the geometric inner-product norm between orthogonal basis generators is exactly unity. All combinatorial symmetries are fully internalized within each pillar's local phase-space calculation (e.g., the rational fractions in $Z_{\mathrm{TOL}}$ and $Z_{\mathrm{MAR}}$). Consequently, inter-pillar coupling coefficients are structurally locked to $s_i = \pm 1$.

\subsubsection{Strict Orthogonality and Pure Partition Factorization}
To combine these pillars without cross-terms ($Z_i Z_j$), the mutual information between them must vanish identically ($I(A;B)=0$). Because the pillars govern strictly orthogonal geometric domains at the zero-flux baseline, their mutual information vanishes, yielding a factorized partition function ($\mathcal{Z}_{\mathrm{total}} = \prod_i \mathcal{Z}_i$). The resulting entropic action $S_\Phi \propto -\ln \mathcal{Z}_{\mathrm{total}} = \sum_i S_i$ is strictly additive. Defining total geometric impedance as this action per unit flux quantum ($Z_{\mathrm{geo}} \equiv S_\Phi / \Phi_{\mathrm{top}}$) maps this directly to $Z_{\mathrm{geo}} = \sum Z_i$, ensuring all higher-order cross-terms vanish.

\subsubsection{Sign Assignment via Variational Impedance}
While the absolute geometric cost of instantiating each component is strictly positive, the observable stable configuration is governed by the \textbf{Variational Impedance}. In this active, closed-loop geometry, direct physical boundary allocations (Capacity, Map, Toll, Margin) represent open-loop structural costs ($s_i = +1$). Conversely, the coordination mechanisms (Protocol, Governor) act as active feedback controllers. By geometrically suppressing uncoordinated phase space across the lattice, they serve as the feedback gain that regulates the open-loop impedance, manifesting as strictly negative reductions ($s_i = -1$).

\vspace{1em}
Therefore:
\begin{equation}
    Z_{\mathrm{geo}} = \sum_{i=1}^6 s_i Z_i .
\end{equation}
\hfill $\blacksquare$

This theorem establishes the unique form of the combination; the derivation of each individual term $Z_i$ from its pillar's geometric constraints is given in \crefrange{subsec:alpha_cap}{subsec:alpha_mar}.


\subsection{The Finite Sum and Causal Coupling}
\label{subsec:finite_sum_numerical}

By the Additive Impedance Theorem, the total geometric impedance is the exact expansion of this sum using the individual structural costs derived in \crefrange{subsec:alpha_cap}{subsec:alpha_mar}:

\begin{equation}\label{eq:alpha_inverse}
\begin{split}
Z_{\text{geo}} = \underbrace{\pi\Delta}_{\text{CAP}}
+ \,\underbrace{\chi}_{\text{MAP}}
- \,\underbrace{\frac{1}{R_M - \sigma}}_{\text{PRO}}
- \,\underbrace{\frac{\chi}{\Delta}}_{\text{GOV}} \\
+ \,\underbrace{\frac{1 \cdot \chi \cdot (R_M - \sigma)}{N^3 \cdot \sigma \cdot R_M}}_{\text{TOL}}
+ \,\underbrace{\frac{1}{L_{\text{embed}} \cdot (\sigma + 1) \cdot \Delta^2}}_{\text{MAR}}
\end{split}
\end{equation}

Summing these geometric components gives the structural index:
\begin{equation}
    Z_{\text{geo}} = \num{\AlphaInvVal}\dots
\end{equation}

Every term is an exact rational combination of the derived geometric bounds except $\pi$, which enters as the universal isotropic ratio $C/\Delta$ (\Cref{subsec:alpha_cap}). The derivation contains strictly zero continuous empirical parameters; precision is therefore limited only by the evaluation of $\pi$.

By the Network Admittance Postulate, the fundamental causal coupling (the probability of non-trivial interaction) of the $E_8$ substrate is exactly the inverse of this impedance:
\begin{equation}
    \kappa = \frac{1}{Z_{geo}} \approx \frac{1}{\num{\AlphaInvVal}}
\end{equation}